\section{Introduction}
\label{sec:intro}

A deployed LLM system makes many small decisions around each answer it writes.
A router decides which model should handle a prompt~\cite{chen2023frugalgpt,ong2024routellm}.
A semantic cache decides whether a stored answer can be served instead of a new one~\cite{bang2023gptcache,schroeder2025vcache}.
An agent decides which tool to call next, or whether to stop and reply~\cite{yao2024taubench}.
Each of these has a small, fixed set of outcomes, yet the usual way to make them is to ask a generative model and parse its text, paying for a full autoregressive call to obtain a few bits.

A recent class of models targets exactly this gap.
These \emph{decision models}, which their vendors call ``System One'' models after the fast mode of dual-process theory~\cite{kahneman2011thinking}, are not autoregressive: given a state and a typed question (pick one of these options, score on this scale, yes or no), they return the answer with a probability for each option in a single pass.
The first commercial one, \jev~\cite{typesafe2026jev}, is priced at a small fraction of a frontier LLM per input token with no output cost, and is advertised as a drop-in replacement for an LLM wherever the output is a decision.
An open-weights reproduction, \kev~\cite{palmer2026kev}, followed within weeks.

Whether this replacement pays is an empirical systems question, and it has two parts that vendor numbers do not separate.
The first is quality: is the cheap decision as good as the expensive one?
The second is structure: what does the decision buy?
A cache guard's decision can remove a generation entirely, so a cheap, accurate guard saves nearly the whole call.
An agent's tool choice sits in front of a generation that happens anyway, because some model must still write the tool's arguments, so a separate decider adds a call instead of removing one.

We study both parts on one harness.
Every decider implements the same interface (a state, a set of options, a returned choice with a confidence), so the same code path scores a frontier LLM, a small hosted LLM, a self-hosted LLM, \jev, \kev, and trained classifiers.
We run them at three decision points on public benchmarks, with repeated calls, clustered confidence intervals, measured dollar cost per call, and measured GPU energy for the self-hosted deciders.

Our findings are as follows.

\textbf{Cache guard: the decision model pays off.}
On GSM-Plus~\cite{li2024gsmplus}, where a question and its nearest cached neighbour often differ by one number, embedding similarity cannot tell a safe reuse from a wrong one, and a dev-tuned threshold reuses almost nothing (66.1\% accuracy).
\jev\ reaches 93.3\%, within 2.6 points of the frontier model, at about 2\% of its price per decision.
A frontier model used as the guard is self-defeating: its decision costs more than the answer it might save.
On two realistic traffic benchmarks \jev\ ties the threshold where similarity works and beats it by 4.4 points where it does not, and a cascade that calls \jev\ only inside the ambiguous similarity band is best or tied for best on all three datasets.

\textbf{Agent control: it does not.}
On \taubench\ airline tasks~\cite{barres2025tau2}, \jev\ alone picks the reference action 81.4\% of the time at $1/22$ of the frontier model's cost and latency.
Inside a real agent loop, however, the arm in which the generator decides and acts in one call is the cheapest of all and matches the best reward.
Splitting the decision out either costs more (the generator asked twice) or loses about five points of reward (\jev\ deciding).

\textbf{Routing: a trained classifier wins.}
On RouterBench~\cite{hu2024routerbench}, no general decider, frontier LLM included, routes usefully from the prompt alone.
A small classifier trained on the cheap model's recorded correctness beats all of them, and even it recovers little.

\textbf{The pattern.}
A standalone decider is worth its call when two conditions hold: the decision avoids a generation, and the evidence it needs is visible in its input.
Cache guarding meets both. Agent control fails the first and routing fails the second.

\textbf{A benchmark flaw.}
While running the cache guard we found that every decider, the frontier model included, stopped at about 71\% accuracy on the two standard semantic-caching benchmarks~\cite{schroeder2025vcache}.
The cause is the labels: they are inferred from classes (transitive closures of pairwise LLM judgements, or ``generated from the same seed'') and not judged on the pair a cache actually meets.
We relabel every pair we use under a written guideline, changing 6.0\% and 15.7\% of labels, and release the corrections with the reasons for each.

We make three contributions:
(1) an independent measurement of the decision-model paradigm at three LLM decision points, in dollars, seconds and joules, against open baselines that run without any API;
(2) a decider interface and harness that make a typed-output decision model interchangeable with chat LLMs and classifiers; and
(3) an analysis of how class-based labelling breaks semantic-cache benchmarks, with corrected labels.
